\subsection{Remaining-Sequence Admission}
\label{sec:admission}

Immediately before each optional stage, \emph{Admission} compares the capture's remaining timeout budget with the estimated duration of the entire remaining draft sequence, including mandatory encoding and saving.
This timing accounts for the budget consumed by queueing and earlier stages.

\smallskip\noindent\textbf{Admission rule.}
The point estimate \(\hat P(\mathcal K_{i,j})\) may underestimate the time needed to finish the remaining sequence.
Admission therefore uses an empirical upper estimate \(U(\mathcal K_{i,j})\) of the remaining processing time for the budget comparison.
\(\mathit{controller}.\textsc{DecideAdmission}\) (Algorithm~\ref{alg:overview}, line~\ref{alg:ov:admission}) admits an optional stage \(k_j\) when
\begin{equation}
  U(\mathcal{K}_{i,j})\leq T_{i,j}
  \quad\text{or}\quad
  \hat P(\mathcal{K}_{i,j})=0.
  \label{eq:admission}
\end{equation}
When \(\hat P(\mathcal K_{i,j})=0\), the predictor has no prior timing measurements for the remaining enabled stages, so the optional stage is admitted to collect initial measurements.
This rule assumes that each configuration can complete an isolated first capture within the Capture Timeout, a capture-level requirement enforced by release validation (Section~\ref{sec:release-process}).

Admission treats \(M\) and \(S\) as separate groups.
Once an \(M\) or \(S\) stage is rejected, admission disables the optional stages in that group.
Until the draft queue drains, admission skips those stages and excludes them from subsequent \(\hat P\) and \(U\) estimates.
This avoids repeatedly toggling the group's visual effects across consecutive draft images.
When the queue drains, both groups are enabled again.

\smallskip\noindent\textbf{Upper-estimate calibration.}
Admission computes the upper estimate by multiplying the point estimate \(\hat P(\mathcal K)\) by an empirical residual factor \(\phi_\tau(\mathcal K)\geq1\):
\begin{equation}
  U(\mathcal{K})
  =\hat P(\mathcal{K})\,\phi_{\tau}(\mathcal{K}).
  \label{eq:upper-estimate}
\end{equation}
After the shared predictor is updated at draft completion, admission calibrates \(\phi_\tau\) for later decisions from observed stage overruns (Algorithm~\ref{alg:overview}, line~\ref{alg:ov:update}).
For each remaining sequence \(\mathcal K_{i,j}\), it forms a measured subsequence \(\mathcal K\) containing only executed stages (\(p_i(k)>0\)).
If \(\mathcal K\) contains a previously measured stage, admission computes a residual factor using the stage-duration predictions \(\hat p_i(k)\) retained before execution:
\begin{equation}
  \phi(\mathcal{K})
  =\frac{\sum_{k\in\mathcal{K}}\max\{\hat p_i(k),p_i(k)\}}
         {\sum_{k\in\mathcal{K}}\hat p_i(k)}.
  \label{eq:residual-factor}
\end{equation}
The per-key maximum preserves each stage's overrun even when another stage finishes faster than predicted, ensuring \(\phi(\mathcal{K})\geq1\).
Normalizing by the total predicted duration expresses stage overruns as a relative correction, making factors comparable across sequences with different predicted durations.

Admission maintains a residual-factor history for each measured sequence.
When a completion yields factors, admission multiplies all existing history weights by \(0.90\) to emphasize recent overruns.
We selected this decay factor at the knee of the trade-off between the overrun rate and the magnitude of \(U\).
Admission then inserts one factor per distinct measured sequence into its history with unit weight.
It selects each history's weighted \(\tau\)-quantile as its residual factor \(\phi_\tau\).
It determines \(\tau\) from the Kish effective sample size \(n_{\mathrm{eff}}\) of the current weights \(w_l\) in that history~\cite{kish-survey-sampling}:
\[
  n_{\mathrm{eff}}
  =\frac{(\sum_l w_l)^2}{\sum_l w_l^2},
  \qquad
  \tau=\frac{n_{\mathrm{eff}}}{n_{\mathrm{eff}}+1}.
\]
This choice of \(\tau\) is motivated by the expected CDF position \(n/(n+1)\) of the maximum of \(n\) i.i.d. continuous samples~\cite{david-order-statistics}, with \(n_{\mathrm{eff}}\) substituting for the sample count.
As the effective sample size grows, the selected quantile moves further into the observed upper tail.
Setting \(\tau\) this way avoids a separately tuned target percentile.
For a sequence without residual history, admission applies the same rule to residual factors pooled across sequences.
If no residual factors have been observed, it sets \(\phi_\tau(\mathcal K)=1\).

\smallskip\noindent\textbf{Watchdog fallback.}
An admitted stage may still run longer than predicted, and it cannot be cancelled once it enters native image processing.
The watchdog therefore triggers fallback if the stage has not completed within the following time budget:
\begin{equation}
  W_{i,j}=\max\{0,T_{i,j}-U(\mathcal{K}_{i,m})\},
  \label{eq:watchdog}
\end{equation}
where \(\mathcal{K}_{i,m}=(k_m)\), so \(U(\mathcal K_{i,m})\) reserves time for mandatory encoding and saving.
Fallback encodes and saves the preserved original input on a separate thread while the admitted stage continues running.
The resources held by the stage are released after it returns.
Queued draft sequences also bypass optional processing until the queue drains, reducing the risk that an overrun propagates to subsequent captures.
